\section{Encoding properties and proofs of exact moments}\label{app:encoding}
We describe the XOR implementation of keySum and checksum and prove the finite identities in Sections~\ref{sec:design} and~\ref{sec:measurement}. We then summarize their exact verification for small parameter values and discuss the AMS observation obtained from the count sum.

\subsection{XOR implementation}\label{app:xor}
In the set model, each side inserts each key at most once. We can therefore accumulate keySum and checksum by bitwise XOR, as in a common implementation of the standard IBLT, and multiply only the count by $\varepsilon(x)$. Under XOR, $-x=x$. Thus, in each difference-table cell, keySum is the XOR of the differing keys that select the cell, and checksum is the XOR of their $\mathrm{chk}$ values. Both are independent of the coding coefficients. In a pure cell, keySum and checksum equal $x_i$ and $\mathrm{chk}(x_i)$ directly, so we check whether checksum equals $\mathrm{chk}(x)$. To delete $x_i$, we subtract $C_v$ from the count and XOR the other two fields once more with $x_i$ and $\mathrm{chk}(x_i)$, respectively. The induction argument for the additive form in the next subsection also applies here. We recover the direction as $C_v\,\varepsilon(x)$, and peeling follows the same process and gives the same result as in the additive form. In a multiset, the same key may be inserted repeatedly. XOR cancels repeated contributions while the count retains them, so multisets require the additive form. If peeling modifies the difference table in place, the receiver can accumulate the two scalars $\sum_vC_v^2$ and $\sum_vC_v$ during cell-by-cell subtraction, without retaining the original difference table.


\subsection{\texorpdfstring{Proof of Proposition~\ref{prop:decode} (ideal decoding equivalence)}{Proof of ideal decoding equivalence}}\label{app:proof-decode}
\begin{proof}
We induct on the number of peeling steps. After each step, the three fields of every cell equal the encoded sum of the remaining keys: the count is $\sum s_i\varepsilon_i$, keySum is $\sum s_i\varepsilon_i x_i$, and checksum is $\sum s_i\varepsilon_i\,\mathrm{chk}(x_i)$. Each sum ranges over the remaining keys that select the cell. This invariant holds initially by~\eqref{eq:fields}.

Suppose that at some step cell $v$ has degree one in the unsigned hypergraph, with unique key $x_i$. Its count is $s_i\varepsilon_i\in\{\pm1\}$, count times keySum equals $x_i$, and checksum equals count times $\mathrm{chk}(x_i)$. The cell hash of $x_i$ also contains $v$, so $v$ passes the pure-cell test. Conversely, a cell of any other degree that passes the test is a false positive, which the ideal-purity assumption excludes. Thus the cells removable under random-sign encoding are exactly the unsigned degree-one cells. To delete $x_i$, we subtract its three-field contribution multiplied by the coding coefficient $s_i\varepsilon_i$ from its $k$ cells. This preserves the invariant and removes the corresponding edge from the unsigned hypergraph. With the same cell-selection order, the two processes reach the same endpoint. Peeling always ends at the 2-core, independently of the selection order, so $F$ depends only on $\mathcal G$.

At termination, we check all fields. Under random signs, a nonempty core can have all-zero counts, so counts alone cannot certify success. The keySum fields can cancel at the same time. Success is misjudged only if all checksums are also zero. This event is included in $E$ in Section~\ref{sec:protocol}.
\end{proof}

\subsection{\texorpdfstring{Proof of Proposition~\ref{prop:moments} (conditional unbiasedness and conditional moments)}{Proof of conditional unbiasedness and conditional moments}}\label{app:proof-moments}
\begin{proof}
Using $H_{ii}=k$ and $\varepsilon_i^2=1$, we obtain the sum of squared counts and the squared sum:
\begin{align*}
\sum_vC_v^2&=\sum_{i,j}\varepsilon_i\varepsilon_jH_{ij}=kd+2\sum_{i<j}\varepsilon_i\varepsilon_jH_{ij},\\
\frac1M\Bigl(\sum_vC_v\Bigr)^2&=\frac{k^2}M\Bigl(\sum_i\varepsilon_i\Bigr)^2=\frac{k^2}Md+2\sum_{i<j}\varepsilon_i\varepsilon_j\frac{k^2}M.
\end{align*}
Dividing the first identity by $k$ gives~\eqref{eq:raw-expand}. Subtracting the second identity from the first gives $kd(1-k/M)+2\sum_{i<j}\varepsilon_i\varepsilon_j(H_{ij}-k^2/M)$. Dividing this by $k(1-k/M)$ gives~\eqref{eq:rs-expand}. For a fixed mapping, $\E[\varepsilon_i\varepsilon_j]=0$ ($i\ne j$), so the conditional expectations of $\widehat d-d$ and $\widetilde d-d$ are both zero.

We now express the coefficients of the cross terms as inner products. Let $w_i\in\mathbb R^M$ be the incidence vector of key $i$ minus its mean $k/M$. Then $\langle w_i,w_j\rangle=H_{ij}-k^2/M$ and $\|w_i\|^2=k(1-k/M)$, so~\eqref{eq:rs-expand} becomes
\begin{equation*}
\widehat d-d=\frac{2}{k(1-k/M)}\sum_{i<j}\langle w_i,w_j\rangle\,\varepsilon_i\varepsilon_j.
\end{equation*}
Given $\mathcal G$, the inner products are constant. All randomness comes from the sign products $\varepsilon_i\varepsilon_j$. The expectation of a product of sign products is nonzero if and only if each subscript occurs an even number of times, in which case the expectation is one.

For the variance, we expand the square and consider two unordered pairs $\{i,j\}$ and $\{i',j'\}$. When the pairs coincide, the product is $\varepsilon_i^2\varepsilon_j^2=1$. When they differ, they share at most one subscript, so some subscript occurs exactly once and the expectation is zero. Thus only the diagonal terms remain:
\begin{equation*}
\Var(\widehat d\mid\mathcal G)=\frac{4}{k^2(1-k/M)^2}\sum_{i<j}\langle w_i,w_j\rangle^2,
\end{equation*}
which gives the equality in~\eqref{eq:condmoments}.

For the third central moment, we expand the cube and consider ordered triples of unordered pairs. If all three pairs coincide, the product is $\varepsilon_i\varepsilon_j$. If exactly two coincide, the product is the sign product of the third pair. Both cases have expectation zero. If the pairs are distinct, every subscript occurs an even number of times exactly when the pairs form a triangle $\{i,j\},\{i,h\},\{j,h\}$. Each triangle corresponds to $3!=6$ ordered triples. Hence
\begin{equation*}
\E[(\widehat d-d)^3\mid\mathcal G]=\frac{8\cdot6}{k^3(1-k/M)^3}\sum_{i<j<h}\Bigl(H_{ij}-\frac{k^2}{M}\Bigr)\Bigl(H_{ih}-\frac{k^2}{M}\Bigr)\Bigl(H_{jh}-\frac{k^2}{M}\Bigr).
\end{equation*}

For the two upper bounds, we use Cauchy--Schwarz: $\langle w_i,w_j\rangle^2\le\|w_i\|^2\|w_j\|^2=[k(1-k/M)]^2$. Substitution into the variance formula shows that each of the $\binom d2$ unordered pairs contributes at most $4$, giving $\Var(\widehat d\mid\mathcal G)\le2d(d-1)$. The component of the vector $\sum_i\varepsilon_iw_i$ on cell $v$ is $C_v-\frac1M\sum_uC_u$, so
\begin{equation*}
\widehat d=\frac{\|\sum_i\varepsilon_iw_i\|^2}{k(1-k/M)}\ge0,
\end{equation*}
and the triangle inequality gives $\|\sum_i\varepsilon_iw_i\|\le\sum_i\|w_i\|=d\sqrt{k(1-k/M)}$, so $\widehat d\le d^2$.

Under the random mapping $\mathcal G$, the conditional mean is always $d$. The law of total variance therefore gives $\Var\widehat d=\E[\Var(\widehat d\mid\mathcal G)]$. The intersection size $H_{ij}$ of two independent uniform $k$-subsets has a hypergeometric distribution with mean $k^2/M$ and variance $k^2(1-k/M)^2/(M-1)$. Substitution into the variance formula shows that each unordered pair contributes $4/(M-1)$, giving $\Var\widehat d=2d(d-1)/(M-1)=s(d)^2$.
\end{proof}

\subsection{Exact verification for small parameter values}\label{app:exact}
We verify the finite identities in Sections~\ref{sec:design} and~\ref{sec:measurement} by exhaustive enumeration for small parameter values. For the six triples $(M,k,d)$ in Table~\ref{tab:exact}, we enumerate all mappings and signs, covering 155,640 states. All arithmetic uses integers and rationals, so the results are exact. In every state, common elements cancel in the difference table. Signed peeling recovers the correct keys and directions, and the unrecovered keys form exactly the unsigned 2-core. For every mapping, $\E[\widehat d\mid\mathcal G]=d$, the conditional variance and third central moment match the formulas above, and $\Var(\widehat d\mid\mathcal G)\le2d(d-1)$. The distribution of $\widehat d$ over the random signs is the same for every direction vector. Table~\ref{tab:exact} summarizes the unconditional and post-failure statistics.

\begin{table}[!htbp]
\centering
\caption{Exact enumeration results for small parameter values.}
\label{tab:exact}
\begin{tabular}{@{}crccccc@{}}
\toprule
$(M,k,d)$ & states & $p_F$ & $\Var\widehat d$ & $\E[\widehat d\mid F]$ & $\Var(\widehat d\mid F)$ & $\Pr(\widehat d=0\mid F)$\\
\midrule
$(4,3,2)$ & 40 & $1/4$ & $4/3$ & 2 & 4 & $1/2$\\
$(4,3,3)$ & 160 & 1 & 4 & 3 & 4 & 0\\
$(4,3,4)$ & 560 & 1 & 8 & 4 & 8 & $27/512$\\
$(5,3,3)$ & 1,760 & $2/5$ & 3 & 3 & 4 & 0\\
$(5,3,4)$ & 11,440 & 1 & 6 & 4 & 6 & $93/8000$\\
$(6,3,4)$ & 141,680 & $2029/4000$ & $24/5$ & 4 & $11256/2029$ & $75/4058$\\
\bottomrule
\end{tabular}
\end{table}

In every row of Table~\ref{tab:exact}, $\Var\widehat d$ equals $s(d)^2=2d(d-1)/(M-1)$ and $\E[\widehat d\mid F]$ equals $d$, consistent with conditional unbiasedness given the mapping in~\eqref{eq:condmoments}. Failure preserves the mean and increases the scale. For the three parameter sets with $p_F<1$, $\Var(\widehat d\mid F)$ is $3$ times, $4/3$ times, and about $1.16$ times the unconditional variance, respectively. When $p_F=1$, failure is certain and the two variances coincide. After failure, $\widehat d$ equals zero with positive probability. For $(4,3,2)$, this probability is $1/2$.

The values in the $(4,3,2)$ row come from a duplicate edge pair. Each key independently and uniformly chooses one of the $\binom43=4$ $3$-subsets. When the two keys choose the same subset, each of these three cells has degree two. There is no pure cell, and peeling stops immediately, so $p_F=1/4$. The counts of the three cells then all equal $\varepsilon_1+\varepsilon_2$, and the remaining cells have zero counts. When the two signs are equal, $\sum_vC_v^2=12$ and $\sum_vC_v=\pm6$, so~\eqref{eq:dhat} gives $\widehat d=4$. When they differ, all counts are zero and $\widehat d=0$. The two cases each have probability $1/2$. Thus $\E[\widehat d\mid F]=2=d$ and $\Var(\widehat d\mid F)=4$. This variance is $3$ times the unconditional variance $4/3$, and $\Pr(\widehat d=0\mid F)=1/2$. This is the smallest instance of the duplicate edge pair structure in Section~\ref{sec:inversion}.


\subsection{An AMS observation independent of failure}\label{app:ams}
Each key enters $k$ cells with the same sign. Thus, under random signs, the normalized count sum equals the sign sum in every sample:
\begin{equation}
\frac1k\sum_vC_v=\sum_{i=1}^d\varepsilon_i.\label{eq:ams}
\end{equation}
The right-hand side depends only on the signs, and its square is the standard AMS second-moment estimator with unit weights\cite{alon1999spacecomplexity}. Rearranging~\eqref{eq:dhat} and~\eqref{eq:raw-energy}, we relate $\widetilde d$ and $\widehat d$ through this AMS term:
\begin{equation}
\widetilde d=\Bigl(1-\frac kM\Bigr)\widehat d+\frac kM\Bigl(\frac1k\sum_vC_v\Bigr)^2.\label{eq:raw-decomp}
\end{equation}

\begin{proposition}[AMS on the same table]\label{prop:ams}
The full distribution of $(\frac1k\sum_vC_v)^2$ is independent of the mapping and hence of ideal failure. Specifically, $\frac1k\sum_vC_v$ takes the value $d-2j$ with probability $2^{-d}\binom dj$, $j=0,\ldots,d$. The squared statistic has mean $d$ and variance $2d(d-1)$, which is $M-1$ times the unconditional variance $s(d)^2$ of $\widehat d$. By~\eqref{eq:raw-decomp}, $\widetilde d$ is a weighted sum of this squared statistic and $\widehat d$, with deterministic weights. Given the mapping, the conditional covariance of the squared statistic with $\widehat d$ is $\frac{4}{k(1-k/M)}\sum_{i<j}(H_{ij}-k^2/M)$.
\end{proposition}

\begin{proof}
The right-hand side of~\eqref{eq:ams} depends only on the signs and is independent of $\mathcal G$. It is therefore independent of every event that depends only on the mapping. $\sum_i\varepsilon_i=d-2j$ if and only if exactly $j$ signs are negative, with probability $2^{-d}\binom dj$. $\E(\sum_i\varepsilon_i)^2=d$, and the fourth moment is $3d^2-2d$, so the variance is $2d(d-1)$. \eqref{eq:raw-decomp} follows by directly rearranging~\eqref{eq:dhat} and~\eqref{eq:raw-energy}. For the covariance formula, we expand both statistics as quadratic forms in the signs and use the fact that sign products for distinct unordered pairs are uncorrelated. We multiply $(\sum_i\varepsilon_i)^2-d=2\sum_{i<j}\varepsilon_i\varepsilon_j$ by~\eqref{eq:rs-expand} and take expectations pair by pair. Averaging over failure gives $2d(d-1)c_F$, which is generally nonzero.
\end{proof}

This observation connects the design to classical AMS. The count array is a random-sign linear sketch whose sum gives exactly the AMS sketch. The AMS sketch is independent of failure, so its distribution remains unchanged however rare failure is. This independence comes at a cost in precision: the AMS estimator has variance $M-1$ times that of $\widehat d$. The estimator $\widehat d$ replaces the square of a single sum with the sum of squares over the $M$ cells, reducing the variance to $2d(d-1)/(M-1)$. Its correlation with failure requires the post-failure analysis. The covariance in Proposition~\ref{prop:ams} shows that the two estimators are generally correlated after failure, so we must allocate risk separately when combining them.

We verify Proposition~\ref{prop:ams} using the exhaustive enumeration procedure in Appendix~\ref{app:exact} for the four parameter sets $(4,3,2)$, $(4,3,4)$, $(5,3,3)$, and $(6,3,4)$, covering 144,040 states. For every mapping, $(\frac1k\sum_vC_v)^2$ has the same distribution, with mean $d$ and variance $2d(d-1)$.
